% ====================================================================
\section{Concluding remarks}\label{sec:concl}
% ====================================================================

\subsection{What the argument rests on}
Apart from classical results (John, Sheppard, Caro--Wei, K\"onig, Blokhuis), every step is proved in this manuscript. The
load-bearing new ingredient is Theorem~\ref{thm:LR}. The geometric induction (Sections~\ref{sec:framework}--\ref{sec:main}) uses
it only as a black box, and all other inputs are elementary. In decreasing order of what we regard as the risk of an error,
the parts that most deserve an expert's attention are the following.
\begin{enumerate}[leftmargin=2em]
\item \emph{Theorem~LR}, in particular the normal form (Lemma~\ref{lem:NF}) and the rescaling at the largest gap by a row
  function, a column function or a constant. Its constants-only and column-only special cases are close to known
  arguments~\cite{Lovett,Rothvoss}. The mixed row/column case rests on Lemmas~\ref{lem:A}, \ref{lem:B} and~\ref{lem:NF}.
\item \emph{Step~A} of Theorem~\ref{thm:Rplus} (Lemma~\ref{lem:stepA}). This is the only place where a density statement (the
  cross-pair count~\eqref{eq:S2}) is converted into a covering statement, at the cost of a fourth root
  (Lemma~\ref{lem:twodist}). An earlier version contained a genuine threshold gap here (fix~F1).
\item The bookkeeping of the Core and Slicing Lemmas, in particular the completeness of the coincidence lists in
  Lemmas~\ref{lem:step0}, \ref{lem:step1} and~\ref{lem:51}, and the verification of (F1)--(F5) in Lemma~\ref{lem:slicing}(d).
\end{enumerate}
If Theorem~\ref{thm:LR} failed, the argument would give only the classical bound $2^{O(d)}$.

\subsection{Computational sanity checks}
Several finite, threshold-free parts of the argument were tested by computer, in the project and by the referees.
We list these tests for completeness. They are \emph{sanity checks only}. Every threshold in the argument
($\smin\ge2^{\Lst}(h+1)$, the cleaning densities, the sides required in the orthogonal splits) is astronomically
larger than any set that can be handled numerically, so no threshold-dependent conclusion was, or could be,
tested. Only inclusions and implications were tested.
\begin{itemize}[leftmargin=1.5em]
\item \emph{Rounding and cleaning} (referee pass on Theorem~LR). John positions were computed numerically (as
  minimum-volume enclosing ellipsoids) on $k$-listable instances, giving $\max|u_x|/s^{1/4}\le1.0002$ (the solver
  tolerance) and $\max|v_w|/(Ks^{1/4})\le1$. Sheppard's formula was compared with Monte Carlo estimates and agreed to
  within $10^{-3}$. The cleaning lemma was tested on adversarial sparse low-rank matrices, and a zero block was always
  found.
\item \emph{Normal form} (same referee pass). An independent implementation of Lemmas~A, B, NF and one round of the LR
  induction was run on $1600$ adversarial block-structured instances: accumulating labels $\kappa\pm2^{-j}$ with $j$ up
  to $40$, ties between row and column functions on blocks, magnitudes $2^{\pm40}$, and up to $5$ types of mixed kinds.
  None of (NF1)--(NF3), the bounds on $|N|$, the identification of the minority set with the support of $H$, or the
  rank bound for $H$ was ever violated.
\item \emph{Genericity in Step~A} (referee pass on Theorem~R$^+$). In $70{,}213$ random grid configurations
  ($n=4,5,6$), $20{,}361$ of them non-generic, no non-generic point fell outside the listed loci.
\item \emph{The case $n=5$} (referee pass on the $n=5$ argument). About $150{,}000$ exact-arithmetic checks of the
  threshold-free inclusions of Steps~0--3 and~5 were run on rainbow-$5$-free families (Johnson-type schemes, products, and
  greedy grid sets), with no violation.
\item \emph{General $n$, author's checks.} $348{,}884$ exact-integer checks for $n=6,7$ and all relevant $(r,p)$ were run,
  covering the non-generic loci of Step~0, the inclusions of Step~1, the characterisation in Lemma~\ref{lem:51}, the
  equivalence behind Lemma~\ref{lem:HR}, and the family conditions in Lemma~\ref{lem:slicing}(d). There was no violation.
  These checks sampled \emph{arbitrary} integer points, not points of rainbow-$n$-free sets (referee item~X6). So they
  do not test the one implication that uses rainbow-$n$-freeness, namely the label lemma behind G1.
\item \emph{General $n$, independent checks} (referee pass on the general-$n$ argument). About $340{,}000$ exact checks
  were run on genuinely rainbow-$n$-free sets for $n=6,7$. The sets were orthogonal unions of integer points on small
  spheres, arranged so that cross distances collide, together with few-distance lattices. The checks included the
  full pipeline of Lemmas~\ref{lem:HR} and~\ref{lem:kill}, the initial family of Theorem~\ref{thm:E1}, slicing rounds with
  actual maximum matchings and K\"onig covers, and Steps~0--2 and~5 of the Core Lemma. They also included the label
  lemma (Lemma~\ref{lem:label}) on label-disjoint generic pairs. There was no violation. On these sets families die
  after one or two slicing rounds, and at most $7$ (for $n=6$) or $12$ (for $n=7$) classes were available. So the
  contradiction in the Core Lemma itself was never exercised; only its steps were.
\end{itemize}

\subsection{Literature}
In the referee pass on Theorem~LR, the following works were searched for prior statements of
Theorem~\ref{thm:LR} or Corollary~\ref{cor:RQlist}:~\cite{HMST}, which proves bounds for nonnegative
integer matrices with bounded average and remarks that binarity is not crucial; Sudakov--Tomon~\cite{SudakovTomon};
Rothvoss~\cite{Rothvoss}; Lovett~\cite{Lovett}; and works citing~\cite{SingerSudan}. No statement of a
$2^{-\tilde O(\sqrt d)}$ bound for $k$-listable real matrices was found. The search was not exhaustive. As stressed in
Remark~\ref{rem:SS}, the possible improvement over the bound stated in~\cite{SingerSudan} for parallel partitions should
be confirmed by experts. Corollary~\ref{cor:RQlist} follows from Rothvoss's argument in a few lines, and it may be regarded
as folklore.

\subsection{Further remarks}
\begin{itemize}[leftmargin=1.5em]
\item The class counts $K_r=12\cdot6^{n-5-r}$, the type counts and all constants are far from optimal. For $n=5$ a referee
  observed that the count in Step~5 already works with $7$ classes instead of $12$, at the cost of larger constants.
\item The super-polynomial factor $2^{\Lst}$ in the thresholds cannot be removed from the class-count statement E1$_r$
  alone. If $Y$ is rainbow-$(n-1)$-free and $o$ is any point, then $Y\cup\{o\}$ is rainbow-$n$-free. A referee observed
  that, for $n=5$, taking $Y$ to be a Johnson-type $5$-distance set and $o$ a suitable weight vector gives a point
  with many more than $12$ classes, each of size polynomial in $D$. So E1 with a threshold polynomial in $D$ would be
  false. With the thresholds used here, E1 is consistent with every example known to us.
\item No super-polynomial lower bound for $f_d(n)$ is known to us, so the gap between $d^{\Omega_n(1)}$ and
  $2^{O_n(\sqrt d\log d)}$ remains wide.
\end{itemize}
